%% sections/03_stm.tex

\section{Analytical State Transition Matrix}
\label{sec:stm}

\subsection{Definition and Variational Equations}
\label{sec:stm_def}

The State Transition Matrix $\STM(t, t_0) \in \mathbb{R}^{6\times6}$
maps perturbations in the initial state to perturbations at time $t$:
\begin{equation}
  \delta\mathbf{x}(t) = \STM(t, t_0)\,\delta\mathbf{x}(t_0),
  \qquad
  \mathbf{x} = (\vq, \vp)^T.
  \label{eq:stm_def}
\end{equation}
It satisfies the variational equation
$\dot{\STM} = \mathbf{A}(t)\,\STM$,
$\STM(t_0,t_0) = \mathbf{I}_6$, with
\begin{equation}
  \mathbf{A}(t) =
  \begin{pmatrix}
    \mathbf{0} & \mathbf{I}_3 \\[4pt]
    \Hess(\vq(t)) & \mathbf{0}
  \end{pmatrix}.
  \label{eq:A_matrix}
\end{equation}

\subsection{Analytical Update within the DKD Sequence}
\label{sec:stm_update}

Partitioning $\STM$ into four $3\times 6$ blocks,
$\STM = (\STM_{qq};\,\STM_{pq})^T$
(where $\STM_{qq} = \partial\vq/\partial\mathbf{x}_0$ and
$\STM_{pq} = \partial\vp/\partial\mathbf{x}_0$), each
sub-step of the DKD sequence propagates $\STM$ analytically:

\paragraph{Drift step $\mathcal{D}(h)$:}
\begin{equation}
  \STM_{qq} \leftarrow \STM_{qq} + h\,\STM_{pq}.
  \label{eq:stm_drift}
\end{equation}

\paragraph{Kick step $\mathcal{K}(h)$:}
\begin{equation}
  \STM_{pq} \leftarrow \STM_{pq} + h\,\Hess(\vq)\,\STM_{qq}.
  \label{eq:stm_kick}
\end{equation}

These updates are derived directly from the linearization of the
Drift and Kick operators and involve the same $\Hess(\vq)$ already
computed for the step-size estimate — no additional force evaluations
are required.

\subsection{Computational Advantage}
\label{sec:stm_cost}

The standard numerical STM approach requires 12 additional
integrations (6 pairs of $\pm\delta$ perturbations in each
phase-space direction). The \AAS\ analytical approach replaces
these with two $3\times6$ matrix-vector products per sub-step
(Equations~\ref{eq:stm_drift}--\ref{eq:stm_kick}), reducing the
STM propagation cost from $O(12N_{\mathrm{steps}})$ to
$O(N_{\mathrm{steps}})$ at negligible per-step overhead.

The state vector is stored as a 42-element vector
$\mathbf{y} \in \mathbb{R}^{42}$: the first 6 elements are
$(\vq, \vp)$ and the remaining 36 are the flattened $\STM$ matrix.
This allows the \AAS\ integrator to serve as a drop-in replacement
for the standard numerical STM propagator in the \astdyn\
orbit determination pipeline \citep{bigi_astdyn_library_2025}.

\subsection{Validation Scope and Required Extensions}
\label{sec:stm_validation_scope}

The current manuscript validates the analytical STM primarily through
the Frobenius-norm mismatch against a finite-difference reference over
a 30-day Apophis arc (Section~\ref{sec:stm_bench}). This demonstrates
internal consistency of the implemented variational update, but it is
not yet a complete validation package for orbit-determination use.

Three extensions are required for a full assessment:
(i) a convergence study versus integration step/precision and
finite-difference perturbation amplitude;
(ii) direct comparison against numerical STM propagation strategies
(central differences and variational-equation integration);
and (iii) end-to-end impact on orbit determination (normal-matrix
conditioning, covariance realism, and post-fit residual statistics).
